\chapter{Intro}
\label{chap:1}
\begin{center}
    Варіант $ = \max \left\{5 ; 10\right\} = 10$
\end{center}
Мета комп'ютерного практикуму: Практична реалізація алгебраїчної атаки на фільтрувальний генератор гами; набуття
навичок роботи з системами комп’ютерної алгебри в SAGE.

\noindent Постановка задачі:
\begin{enumerate}
    \item Знайти функції мінімального степеня ідеалів $\langle f \rangle$ та $\langle f \oplus 1\rangle$ за допомогою побудови
        базису Грьобнера. Якщо побудова базису для одного з ідеалів $\langle f \rangle$ або $\langle f \oplus 1\rangle$ є занадто
        трудомісткою з точки зору обчислювальних ресурсів, то будувати лише один базис -- за умови, що цього буде достатньо для 
        проведення атаки.
    \item Визначити кількість рівнянь, необхідних для відновлення початкового стану (якщо не вдається визначити аналітично, 
        то емпірични шляхом підбирати). Побудувати систему рівнянь меншого степеня відносно початкового стану генератора.
    \item Знайти розв’язки отриманої системи рівнянь. Зауважимо, що початковий стан за умовою комп’ютерного практикуму 
        є ненульовим вектором.
    \item Перевірити, що початковий стан відновлено правильно, шляхом генерування відрізка гами відповідної довжини й 
        порівнявши його з вхідними даними.
\end{enumerate}

\section{Вхідні дані}
\begin{enumerate}
    \item Довжина регістру 64 біти. Поліном зворотного зв’язку є наступним:
        \begin{align*}
            p(x) = & x^{64} \oplus x^{63} \oplus x^{62} \oplus x^{60} \oplus x^{59} \oplus x^{58} \oplus x^{57} \oplus x^{56} \oplus x^{53} \oplus x^{50} \oplus x^{47} \oplus x^{45} \oplus x^{44} \oplus x^{43} \oplus \\
            & \oplus x^{42} \oplus x^{41} \oplus x^{40} \oplus x^{39} \oplus x^{38}\oplus x^{37} \oplus x^{36} \oplus x^{34} \oplus x^{32} \oplus x^{30} \oplus x^{28} \oplus x^{24} \oplus x^{18} \oplus \\
            & \oplus x^{15} \oplus x^{14} \oplus x^{13} \oplus x^{11} \oplus x^{9}\oplus x^{6} \oplus x^{4} \oplus 1.
        \end{align*}
    \item Фільтрувальні функції $f(x_{0}, \ldots, x_{63})$ генератора гами збережені у \textbf{filterfuncV10.txt}.
    \item Відрізом гами $\gamma_{1}, \ldots, \gamma_{N}$ довжини 50000 бітів збережений у файлі \textbf{gammaV10.txt}.
\end{enumerate}

\section{Хід виконання роботи}

Я встановив Sage в свій WSL 2 ubuntu за інструкцією~\url{https://doc.sagemath.org/html/en/installation/index.html}.


\section{Coding section}

\section{Висновки}
